\section{Density of the even tests}\label{sec:density}

The Hessian identity identifies affine functions as the possible
obstructions to density.  We now show that these are the only
obstructions, using cutoffs in the $x$ coordinates.

\begin{lemma}[Density of even tests]\label{lem:density}
Let $d\mu=\rho(x)\,dx$ be a probability measure on $\R^n$ with a smooth,
strictly positive, even density.  Let $\tau$ be a smooth even symmetric
matrix field with $0<\tau(x)\le\Lambda I$ for some finite $\Lambda$, and set
\[
 Af=x\cdot\nabla f-\tau:D^2f.
\]
Suppose that $A$ has the integration formula
\[
 \E[hAf]=\E[\nabla h\cdot\tau\nabla f]
 \qquad(f,h\in C_c^\infty(\R^n))
\]
and satisfies the compact Hessian identity~\eqref{eq:hessian}.
Then
\[
 \overline{\{A(A-1)f:f\in C_c^\infty(\R^n),\ f\text{ even}\}}^{\,L^2(\mu)}
 =\{g\in L^2(\mu):g\text{ even},\ \E g=0\}.
\]
\end{lemma}

\begin{proof}
Write $\Gamma(f,h)=\nabla f\cdot\tau\nabla h$ and
$\Gamma(f)=\Gamma(f,f)$.  The integration formula also holds when one
factor is smooth and the other has compact support, by inserting a
cutoff around that support.  In particular $\E Ak=0$ for compactly
supported smooth $k$.  Since $A$ preserves parity, the displayed range
consists of centered even functions.

Let $g$ be a centered even $L^2(\mu)$ function orthogonal to this range.
It also annihilates $A(A-1)f$ for odd $f$, because that image is odd.
Formal symmetry therefore gives
\[
 \mathcal P g=0\quad\text{in distributions},\qquad \mathcal P=A(A-1).
\]
Indeed, if $\mathcal P^t$ denotes the transpose with respect to
Lebesgue measure and $d\mu=\rho\,dx$, weighted formal symmetry gives
$\mathcal P^t(\rho g)=\rho\mathcal P g$ in distributions.
Orthogonality makes the left side zero, and smooth positivity of
$\rho$ gives the displayed equation.  The principal symbol of $\mathcal P$ is
$(\xi\cdot\tau(x)\xi)^2$, which is positive for every $\xi\ne0$.
Interior elliptic regularity consequently gives $g\in C^\infty$
\cite[Theorem~14.2]{Dyatlov2026}.  Only positivity on compact sets is
needed here.

\smallskip
\noindent\emph{Splitting the equation in $L^2$.}
Choose $0\le\eta_r\le1$ smooth, equal to one on $B_r$, zero outside
$B_{2r}$, with $|\nabla\eta_r|\le C/r$ and $|D^2\eta_r|\le C/r^2$,
where $r\ge1$.  The upper bound on $\tau$ gives constants $B,K$,
independent of $r$, such that
\begin{equation}\label{eq:density-cutoff}
 \Gamma(\eta_r)^{1/2}\le B/r,
 \qquad |A\eta_r|\le K.
\end{equation}
For the second estimate, use $|x|\le2r$ on the support of the cutoff
derivatives.  Those derivatives, including $A\eta_r$, vanish outside
$B_{2r}\setminus B_r$.

Put $v=Ag$, initially just a smooth function, and write
\[
 G=\|g\|_2,\qquad
 X_r=\|\eta_r^2v\|_2,\qquad
 Y_r^2=\E[\eta_r^2\Gamma(g)].
\]
Here $G$ is finite by hypothesis, while $X_r,Y_r$ are finite by compact
support.
Testing $\mathcal P g=0$ against $\eta_r^4g$ gives
$\E[v(A-1)(\eta_r^4g)]=0$.
The product rule $A(ab)=aAb+bAa-2\Gamma(a,b)$ yields the exact expansion
\[
 \begin{split}
 (A-1)(\eta_r^4g)
 ={}&\eta_r^4v
 +g\bigl(4\eta_r^3A\eta_r-12\eta_r^2\Gamma(\eta_r)-\eta_r^4\bigr)\\
 &-8\eta_r^3\Gamma(\eta_r,g).
 \end{split}
\]
Thus Cauchy--Schwarz and~\eqref{eq:density-cutoff} imply
\begin{equation}\label{eq:density-two-estimates}
 X_r\le aG+\frac{8B}{r}Y_r,
 \qquad
 Y_r^2\le GX_r+\frac{2B}{r}GY_r,
 \qquad a=1+4K+12B^2.
\end{equation}
The first inequality is immediate if $X_r=0$; otherwise divide the
resulting estimate for $X_r^2$ by $X_r$.  For the second, integrate by
parts in $\E[\eta_r^2gAg]$:
\[
 Y_r^2=\E[\eta_r^2gv]-2\E[\eta_rg\Gamma(\eta_r,g)].
\]
Combining~\eqref{eq:density-two-estimates} gives
$Y_r^2\le aG^2+10BGY_r$, so both $Y_r$ and $X_r$ are bounded uniformly
in $r$.  Exhaustion by balls proves $v=Ag\in L^2(\mu)$.  We may now set
\[
 w=g-v\in L^2(\mu),\qquad Aw=0,\qquad Av=v.
\]

\smallskip
\noindent\emph{The harmonic part is constant.}
Testing $Aw=0$ with $\eta_r^2w$ and using
\eqref{eq:density-cutoff} gives
\[
 \bigl(\E[\eta_r^2\Gamma(w)]\bigr)^{1/2}
 \le\frac{2B}{r}\|w\|_2.
\]
On any fixed ball, let $r\to\infty$.  Then $\Gamma(w)=0$ there;
strict positivity of $\tau$ shows that $w$ is constant.

\smallskip
\noindent\emph{The eigenvalue-one part is linear.}
The analogous test of $Av=v$ gives, with
$Z_r^2=\E[\eta_r^2\Gamma(v)]$,
\[
 Z_r^2\le\|v\|_2^2+\frac{2B}{r}\|v\|_2Z_r.
\]
Hence $\E\Gamma(v)<\infty$.  For $f_r=\eta_rv\in C_c^\infty$,
the product rule now proves directly that
\begin{equation}\label{eq:density-approximation}
 f_r\longrightarrow v,\qquad Af_r\longrightarrow v
 \quad\text{in }L^2(\mu).
\end{equation}
Indeed,
\[
 Af_r=\eta_rv+vA\eta_r-2\Gamma(\eta_r,v).
\]
The middle term has norm at most
$K\|v\mathbf{1}_{B_{2r}\setminus B_r}\|_2\to0$.  For the last term, use
\[
 \|\Gamma(\eta_r,v)\|_2
 \le\frac{B}{r}(\E\Gamma(v))^{1/2}\longrightarrow0.
\]
Apply~\eqref{eq:hessian} to these compactly supported functions:
\[
 \E\bigl[\|\tau^{1/2}D^2f_r\,\tau^{1/2}\|_{\mathrm{HS}}^2\bigr]
 =\|Af_r\|_2^2-\E[f_rAf_r]\longrightarrow0.
\]
On each fixed ball, $f_r=v$ for all sufficiently large $r$.  Positivity
of the density and of $\tau$ therefore forces $D^2v=0$ on that ball.
Thus $v$ is affine, and $Av=v$ removes its constant term.

We have proved that every $L^2$ distributional solution of $\mathcal P g=0$ is
affine.  Our $g=w+v$ is even, so it is constant, and its mean is zero.
Hence $g=0$.  Taking the orthogonal complement in the centered even
subspace proves the lemma.
\end{proof}
